
% ================================================================
% CHAPTER 7: CONCLUSION
% ================================================================

\section{Summary of Contributions}

The starting point for this project was a paper on arXiv (2511.10375) describing
a knowledge-graph RAG pipeline called TruthfulRAG v4. The goal here was
not just to reproduce that system but to push it further: to fix the gaps that
even v4 left open.

The result is TruthfulRAG v5, a fully local, domain-agnostic, conflict-resolving
question-answering pipeline. The core contributions are:

\begin{enumerate}\itemsep0pt\parsep0pt
    \item \textbf{Faithful v4 re-implementation}: The three-module architecture
    from arXiv:2511.10375 is reproduced in Python using open-source tools,
    giving the academic community a testable reference implementation.

    \item \textbf{Nine algorithmic extensions}:
    \begin{enumerate}\itemsep0pt\parsep0pt
        \item[\textbf{C1}] \textbf{Temporal decay}: $e^{-\lambda t}$ with
        $\lambda = 0.08$, giving an 8.7-year knowledge half-life applied at query time.
        \item[\textbf{C2}] \textbf{Year-per-triple storage}: Every graph edge
        carries the year the underlying fact was published.
        \item[\textbf{C3}] \textbf{Cross-document corroboration}: A
        logarithmic bonus ($1 + \ln(1+n) \times 0.8$) rewards facts attested
        by multiple independent sources.
        \item[\textbf{C4}] \textbf{Hybrid BM25 + semantic retrieval (RRF)}:
        Keyword and dense-embedding rankings are fused via Reciprocal Rank Fusion.
        \item[\textbf{C5}] \textbf{Gap-threshold conflict elimination}:
        When the score gap between conflicting facts exceeds 0.30, the weaker
        fact is removed before the LLM sees the prompt (gap: 0.028 in v4 vs 1.126 in v5).
        \item[\textbf{C6}] \textbf{Adaptive entropy skip}: Between 24\,\% and
        40\,\% of LLM calls are skipped when parametric entropy is already low
        (measured across all four test domains).
        \item[\textbf{C7}] \textbf{Calibrated confidence score}: A
        three-component score (entropy reduction, source support, temporal
        recency) replaces the fixed-100\,\% certainty of standard RAG.
        \item[\textbf{C8}] \textbf{Explanation audit trail}: Every conflict
        resolution decision is logged with the specific numerical reason.
        \item[\textbf{C9}] \textbf{Automatic schema inference}: An LLM call
        discovers entity and relation vocabulary from the corpus itself.
    \end{enumerate}

    \item \textbf{Production hardening}: Apostrophe injection in Cypher
    strings, malformed LLM JSON, Windows UTF-8 mismatches, APOC absence,
    and corpus dropdown state issues are all handled gracefully.
\end{enumerate}

% ----------------------------------------------------------------
\section{Directions for Further Work}

Four directions stand out as the highest-priority extensions:

\begin{enumerate}\itemsep0pt\parsep0pt
    \item \textbf{Multi-hop graph reasoning.} The retriever currently operates
    on single edges. BFS or PPR path-following retrieval would handle two-edge
    questions at the cost of higher query latency.

    \item \textbf{Domain-adaptive decay rate.} The $\lambda = 0.08$ half-life
    is reasonable for general knowledge but wrong for medical guidelines or
    constitutional law. Auto-selecting $\lambda$ from the inferred domain
    class would make the decay meaningful.

    \item \textbf{Persistent, accumulative graph.} Currently, loading a new
    corpus clears Neo4j entirely. A better design would timestamp old facts
    rather than deleting them, enabling long-term knowledge-base maintenance.

    \item \textbf{Cross-encoder re-ranking and multilingual support.} A
    cross-encoder after PPR scoring would surface the best evidence regardless
    of graph position. A multilingual sentence encoder would extend the system
    to Hindi, Meitei, Tamil, and other regional languages without
    architectural change.
\end{enumerate}

% ----------------------------------------------------------------
\section{Closing Remarks}

Standard RAG systems present every answer with the same implicit confidence:
maximum. That is dishonest. When evidence is contradictory, or retrieved facts
are thirty years old, a system that still says ``here is the answer'' without
qualification is not helpful, it is misleading.

TruthfulRAG v5 takes the opposite position: honesty is a feature, not a
limitation. A 42\,\% confidence score in a high-stakes domain is more useful
than a system that quietly swallows its own uncertainty.

Perhaps the more unexpected finding is that building such a system does not
require industry-scale resources. This was built by one student, on a laptop,
using freely downloadable tools, without sending a single document to an
external server. The organisations that most need trustworthy AI --- public
hospitals, lower courts, government archives --- are exactly the ones that cannot
afford cloud AI subscriptions. A system like this could run for them today,
on hardware they already own.
